\FifthArticle{11}
Various methods can be used to determine numerically the roots of algebraic equations consisting of only three terms. Here these methods admit an elegance and convenience far beyond the laborious operations unavoidable for equations of a less simple form. Such methods therefore deserve a separate exposition, especially since equations of this form occur often enough.

This applies first to the expansion of the roots into infinite series. In fact, \textit{every} root of an equation with three terms, whether real or imaginary, can be expressed by a convergent series with a simple law of progression. For several reasons, however, I shall exclude \textit{this} mode of solution entirely from the present discussion. I only observe here that the rate of convergence depends on the relation between the coefficients: the closer this relation comes to that for which the equation has two equal roots, the slower the convergence; and in this limiting case itself it is slower than that of any decreasing geometric progression. However noteworthy these series may be from a general theoretical point of view, in practice one will always prefer the \textit{indirect} methods to be developed in the following articles, except when the series converge very rapidly.

\FifthArticle{12}
To find the \textit{real} roots I use my auxiliary table for logarithms of sums and differences, first printed in 1810, or, when greater accuracy is required than logarithms of five digits can give,\SourceLeaf{093} the similar but extended table of \textsc{Matthiessen}. I have already published a few special applications of this procedure: to the solution of quadratic equations, in the twentieth edition of \textsc{Vega}'s logarithmic handbook, published in 1840; and to the solution of the cubic equation used to determine the true anomaly in parabolic motion, in No. 474 of the Astronomische Nachrichten. In the latter place the general applicability of the procedure to all algebraic equations with three terms was also pointed out. Although the treatment of this entirely elementary subject presents no difficulty, the considerable variety of cases will make a clear division of them and a collection of rules ready for use welcome enough to justify a few pages.

Instead of these auxiliary logarithmic tables one can also use the ordinary tables of trigonometric logarithms. The former, however, are generally more convenient for the present purpose and give twice the accuracy of the latter. I would therefore restrict the use of trigonometric tables for this task to the rare case in which one wishes to exceed even the accuracy attainable with logarithms of seven digits, and can use for this purpose the well-known logarithms of ten digits in \textsc{Vlacq}'s or \textsc{Vega}'s Thesaurus. Moreover, when auxiliary logarithms are used, twice as many cases must be distinguished as when trigonometric logarithms are used. This must not be regarded as a disadvantage: once the complete general classification is available, it is easy to assign each particular case to its class, and the actual indirect task is much easier when the whole class has only half the extent. But for precisely this reason the general classification is shorter and more convenient to present for solution by trigonometric logarithms. I shall therefore give it first, since the classification for the other mode of solution then follows of itself.

\FifthArticle{13}
The method is applied directly only to determining the \textit{positive} roots of a given equation. The negative roots\SourceLeaf{094} are obtained by applying the same procedure to the equation resulting from the given one on introducing the quantity opposite to the original unknown.

I put the equation in the form
\[
x^{m+n}\pm e x^m\pm f=0
\]
where $m$, $n$, $e$, $f$ denote given \textit{positive} quantities. Depending on the combination of signs, this form actually comprises four different cases; but the first, in which the upper signs are chosen both times, is excluded, since the equation
\[
x^{m+n}+e x^m+f=0
\]
clearly cannot have a positive root. Moreover, it is permissible to assume that $m$ and $n$ (which are understood to be integers, although the applicability of the method itself is independent of this) have no common divisor, since every other case is easily reduced to this one. Finally, for brevity I shall write
\[
\frac{f^n}{e^{m+n}}=\lambda
\]

\Needspace{6\baselineskip}\begin{center}\textit{First Form.}\end{center}
\[
x^{m+n}+e x^m-f=0
\]
Introduce an angle $\theta$, always to be taken in the first quadrant, so that
\[
\frac{x^{m+n}}{f}=\sin\theta^2,\quad
\frac{e x^m}{f}=\cos\theta^2
\]
and hence (I)
\[
x^{m+n}=f\sin\theta^2,\quad
x^m=\frac{f\cos\theta^2}{e},\quad
x^n=e\mathrm{tang}\theta^2
\]
Eliminating $x$ gives the equation
\[
\lambda=\frac{\sin\theta^{2m}}{\cos\theta^{2m+2n}}
\]
from which $\theta$ must be determined. It is easy to see that the right-hand side of this equation, regarded as a function of an indeterminate quantity $\theta$, increases from $0$\SourceLeaf{095} to $\infty$ as $\theta$ passes through all values from $0$ to $90^0$, and that there is therefore one, and only one, value of $\theta$ satisfying the equation. Once it has been found, $x$ is obtained from one of the formulas I. Notice that $\theta=45^0$ when $\lambda=2^n$, and consequently $\theta$ is to be sought in the first octant when $\lambda$ is smaller, and in the second when $\lambda$ is greater than $2^n$.

\Needspace{6\baselineskip}\begin{center}\textit{Second Form.}\end{center}
\[
x^{m+n}-ex^m-f=0
\]
Here put
\[
fx^{-m-n}=\sin\theta^2,\quad ex^{-n}=\cos\theta^2
\]
or (I)
\[
x^{m+n}=\frac{f}{\sin\theta^2},\quad
x^n=\frac{e}{\cos\theta^2},\quad
x^m=\frac{f\mathrm{cotang}\theta^2}{e}
\]
Thus $\theta$ is to be determined from the equation
\[
\lambda=\frac{\sin\theta^{2n}}{\cos\theta^{2m+2n}}
\]
which can be done in one and only one way. The value of $x$ is then found from one of the equations I. The angle $\theta$ lies in the first or second octant according as $\lambda$ is smaller or greater than $2^m$.

\Needspace{6\baselineskip}\begin{center}\textit{Third Form.}\end{center}
\[
x^{m+n}-ex^m+f=0
\]
Here put
\[
\frac{x^n}{e}=\sin\theta^2,\quad
\frac{fx^{-m}}{e}=\cos\theta^2
\]
or (I)
\[
x^{m+n}=f\mathrm{tang}\theta^2,\quad
x^m=\frac{f}{e\cos\theta^2},\quad
x^n=e\sin\theta^2
\]
One of these formulas will serve to determine $x$ as soon as the value of $\theta$ has been found. This is obtained by solving the equation
\[
\lambda=\cos\theta^{2n}\sin\theta^{2m}
\]
Since the right-hand side of this equation, regarded as a function of an indeterminate quantity $\theta$, vanishes both for $\theta=0$ and for $\theta=90^0$, there must be a greatest value between them. Since the differential of the\SourceLeaf{096} logarithm of this function $=(2m\mathrm{cotg}\theta-2n\mathrm{tang}\theta)\,\mathrm{d}\theta$, the greatest value occurs for $\theta=\theta^*$, on putting $\surd\frac{m}{n}=\mathrm{tang}\theta^*$. The function therefore increases from $0$ to its greatest value, evidently
\[
=\frac{m^m n^n}{(m+n)^{m+n}}
\]
and then decreases to $0$, as $\theta$ increases from $0$ to $\theta^*$ and thence to $90^0$. Thus the maximum value is always greater than the value for $\theta=45^0$, that is, greater than $\frac{1}{2^{m+n}}$, except when $m=n$, in which case $\frac{1}{2^{m+n}}$ itself is the maximum value.

It follows that, according as $\lambda$ is greater than
\[
\frac{m^m n^n}{(m+n)^{m+n}}
\]
or smaller, the equation $\lambda=\cos\theta^{2n}\sin\theta^{2m}$ cannot be satisfied at all, or can be satisfied by two distinct values of $\theta$. In the former case the equation $x^{m+n}-ex^m+f=0$ has no (positive) root; in the latter it has two. In the special case where
\[
\lambda=\frac{m^m n^n}{(m+n)^{m+n}}
\]
the two solutions coincide and the equation has two equal roots, for which any of the three formulas may be used at will:
\[
x^{m+n}=\frac{fm}{n},\quad
x^m=\frac{f(m+n)}{en},\quad
x^n=\frac{em}{m+n}
\]

As for the octants in which the values of $\theta$ fall when two solutions actually exist, it is easy to see that, if $\lambda$ is greater than $\frac{1}{2^{m+n}}$, both values of $\theta$ lie in the same octant as $\theta^*$, namely in the first or second according as $m$ is smaller or greater than $n$. If, however, $\lambda$ is smaller than $\frac{1}{2^{m+n}}$, one value of $\theta$ is to be sought in the first octant and the other in the second. In the special case where $\lambda=\frac{1}{2^{m+n}}$, $45^0$ itself is one value of $\theta$, and the other lies in the same octant as $\theta^*$.

One may also note the consequence, easily obtained from this analysis of all three forms, that our equation (insofar as we assume that $m$ and $n$ have no common divisor) cannot have more than three real roots, as is also known on other grounds.
\SourceLeaf{097}

